% Generated draft: requires mathematical and source review.
\documentclass[11pt]{article}
\usepackage[margin=25mm]{geometry}
\usepackage{fontspec}
\setmainfont{Latin Modern Roman}
\usepackage{amsmath,amssymb,mathtools,unicode-math}
\setmathfont{Latin Modern Math}
\usepackage{enumitem,needspace,xurl,hyperref}
\setlength{\parindent}{0pt}
\setlength{\parskip}{5pt}
\setlength{\emergencystretch}{3em}
\newcommand{\N}{\mathbb{N}}
\newcommand{\Z}{\mathbb{Z}}
\newcommand{\Q}{\mathbb{Q}}
\newcommand{\R}{\mathbb{R}}
\newcommand{\C}{\mathbb{C}}
\newcommand{\F}{\mathbb{F}}
\newcommand{\A}{\mathbb{A}}
\newcommand{\PP}{\mathbb{P}}
\newcommand{\E}{\mathbb{E}}
\newcommand{\eps}{\varepsilon}
\newcommand{\dd}{\,\mathrm{d}}
\begin{document}
\section*{UnsolvedMath ID: 800146}
\textbf{Source problem code:} catalogue entry
\subsection{Definitions and mathematical statement}
For an integer \(n\ge 2\), let \(P^+(n)\) be the largest prime divisor of \(n\). Define the real Dirichlet character modulo \(4\) by
\[
\chi_{-4}(m)=\begin{cases}
0,&m\equiv 0,2\pmod 4,\\
1,&m\equiv 1\pmod 4,\\
-1,&m\equiv 3\pmod 4.
\end{cases}
\]
For real \(x\ge 2\), define the explicitly truncated sum
\[
S(x)=\sum_{\substack{n\in\Z\\2\le n\le x}}\chi_{-4}\bigl(P^+(n)\bigr).
\]
Thus \(S(x)\) is the number of integers in this range whose largest prime divisor is congruent to \(1\pmod 4\), minus the number whose largest prime divisor is congruent to \(3\pmod 4\).

The supplied source reports numerical evidence of a bias toward the latter class and describes the asymptotic behavior of its sum as unknown. These are supplied source assessments, not independently verified conclusions.

The source asks to prove or disprove the displayed inequality \(\sum_{n\le x}\chi_{-4}(P^+(n))<0\), followed by the quantifier ``for all \(n\ge 2\).'' This quantifier does not bind the parameter \(x\). A possible intended statement, requiring confirmation, is
\[
S(x)<0\qquad\text{for every real }x\ge 2.
\]
For the defined sum, quantifying over integer cutoffs \(x\ge 2\) gives an equivalent statement. The strict inequality and starting cutoff above reproduce the apparent intended reading without asserting its correctness. Moreover, the source does not specify how its original sum treats \(n=1\), for which the stated largest-prime-factor definition is unavailable. Omitting that term, as in \(S(x)\), is an explicit provisional interpretation requiring review.

The source additionally asks, if possible, for an asymptotic formula for the sum as \(x\to\infty\). No proposed main term, error bound, or additional hypothesis is supplied.
\subsection{Short English statement}
The question concerns whether integers whose largest prime divisor is \(3\pmod 4\) always outnumber those whose largest prime divisor is \(1\pmod 4\) up to each cutoff, and optionally asks for an asymptotic formula for the difference. The cutoff quantifier and treatment of \(n=1\) require clarification; the strict inequality and initial cutoff are retained provisionally.
\subsection{Sources}
\begin{itemize}
\item Supplied local UnsolvedMath registry record: permanent ID \(800146\), exact source code \texttt{catalogue entry}, title ``Comparative Prime Number Theory --- Problem 15.'' The registry labels the problem open and lists Hamieh, Kadiri, Martin, and Ng as proposers, with proposed year \(2024\). These are supplied registry metadata, not independently verified mathematical or literature assessments.
\item Hamieh, Kadiri, Martin, and Ng, \emph{Comparative Prime Number Theory Problem List} (2024), Problem (\#15), \url{https://arxiv.org/abs/2407.03530}. The supplied provenance records access on 2026-10-16 and extraction by \texttt{arxiv\_latex}; this access information is supplied provenance, not a claim of independent consultation.
\item The supplied problem text credits Alexandre Bailleul, Universit\'e Paris--Saclay, and states that the question was raised by Nilotpal Sinha, using the reference label [Sinha]. No bibliographic entry or link resolving [Sinha] is supplied.
\end{itemize}
\end{document}
